\documentclass[conference]{IEEEtran}
\IEEEoverridecommandlockouts
% The preceding line is only needed to identify funding in the first footnote. If that is unneeded, please comment it out.
\usepackage{float}
\usepackage[table]{xcolor}
\usepackage{subcaption}
\usepackage{multirow}
\usepackage{graphicx}
\usepackage{amsmath,amssymb,amsfonts}
\usepackage{tabularx}
\usepackage{booktabs}
\usepackage{microtype}
\usepackage{parskip}

\begin{document}

\title{P2.2 - Users' perception analysis of cyberattack consequences using NLP*\\ 

\footnotesize \textsuperscript{*} https://www.overleaf.com/read/tdbvzhkytcww\#3670b1}

\author{%
Ebadullah Wardak, Kaiming Meng, Lawrence Chen, Oskar Aria Attar\\
CS273\\
https://github.com/RAILGUN1124/CS273-HW2.2/tree/main \\ 
San Jose State University
}

\maketitle

\begin{abstract}

\end{abstract}


\section{Introduction}

\cite{abri2021explainableusers}
\section{Datasets}






\section{Human-Based Similarity Matrix}

\subsection{Preprocessing}

The card-sorting dataset used in this study was collected from a user study in which participants grouped cyber-attack consequences based on perceived similarity. The dataset is stored in \texttt{all\_data.xlsx}. Each row represents a participant's grouping vector, while each column corresponds to a cyber-attack consequence description.

The dataset is structured as follows:

\begin{itemize}
\item Column A: Participant identifier
\item Column B: Group identifier created by the participant
\item Columns C--AZ: Binary document vectors representing the 50 cyber-attack consequences
\end{itemize}

Each binary value indicates whether a specific consequence belongs to a participant-defined group. During preprocessing, only the document vector columns (C--AZ) were extracted. Columns containing only missing values were removed, and empty column labels were discarded to ensure consistent indexing of the consequences.

The resulting matrix represents the relationship between participants and consequences:

\[
\text{Participants} \times \text{Consequences}
\]

This matrix captures how users grouped cyber-attack consequences according to their perception.

\subsection{Method}

To quantify human perception of similarity between cyber-attack consequences, a pairwise similarity matrix was computed using the Spearman rank correlation coefficient. This approach follows the methodology described by Abri et al.~\cite{abri2021explainableusers}, where the similarity between two consequences is determined based on the consistency with which participants grouped them together.

Let $C_i$ and $C_j$ denote the participant grouping vectors for consequences $i$ and $j$, respectively. The similarity between the two consequences is defined as

\[
\rho_{ij} = \text{Spearman}(C_i, C_j)
\]

where $\rho_{ij}$ ranges between $-1$ and $1$. A value of $1$ indicates that two consequences are consistently grouped together by participants, while a value near $0$ indicates no relationship. Negative values indicate opposing grouping patterns.

The resulting matrix is a $50 \times 50$ similarity matrix where both rows and columns correspond to cyber-attack consequences. Each element represents the Spearman correlation between the grouping vectors of two consequences. This matrix serves as the human-based similarity representation of user perception.

\subsection{Results}

The resulting human-based similarity matrix is visualized as a heatmap in Figure~\ref{fig:human_similarity}. The diagonal elements of the matrix have values equal to $1$, indicating perfect correlation between each consequence and itself.

\begin{figure}[h]
\centering
\includegraphics[width=0.9\linewidth]{step1/human_similarity_matrix.png}
\caption{Human-based similarity matrix computed using Spearman rank correlation between participant card-sorting vectors.}
\label{fig:human_similarity}
\end{figure}

Several regions of the heatmap exhibit clusters of higher similarity values, indicating groups of consequences that participants frequently categorized together. However, the matrix also contains many values close to zero, suggesting weak similarity between many consequence pairs. This pattern reflects the variability in how participants conceptualized cyber-attack consequences.

The human similarity matrix serves as the ground truth representation of user perception. In subsequent steps, this matrix will be compared with similarity matrices generated by various NLP techniques, including Bag-of-Words, TF-IDF, Latent Semantic Analysis, and WordNet.
\cite{abri2021explainableusers}.


\section{NLP Vectorization}



\section{WordNet-Based Similarity Analysis}

\subsection{Preprocessing}

To compute semantic similarity between the 50 cyber-attack consequence descriptions, several preprocessing steps were applied before extracting WordNet synsets.

First, each description was tokenized and converted to lowercase. Non-alphabetic tokens such as punctuation and numbers were removed to ensure that only meaningful lexical tokens were considered.

Next, Part-of-Speech (POS) tagging was applied using the Penn Treebank tagset. Only words belonging to the four WordNet-supported categories were retained: nouns, verbs, adjectives, and adverbs.

Four preprocessing configurations were evaluated:

\begin{table}[h]
\centering
\begin{tabular}{ccc}
\hline
Configuration & Stopword Removal & Lemmatization \\
\hline
cfg01 & No & No \\
cfg02 & No & Yes \\
cfg03 & Yes & No \\
cfg04 & Yes & Yes \\
\hline
\end{tabular}
\caption{WordNet preprocessing configurations.}
\label{tab:WNbased_configs}
\end{table}

Stopwords were removed using the NLTK English stopword list. When lemmatization was enabled, words were reduced to their base form using the WordNet lemmatizer to ensure proper mapping to WordNet synsets.

For each remaining token, the most frequent WordNet synset was selected.

\subsection{Method}

Semantic similarity between two attack descriptions was computed using a WordNet-based Greedy Lemma Aligning Overlap (GLAO) approach.

Let $S_1$ and $S_2$ denote the synset sets extracted from two descriptions. For each synset in $S_1$, the algorithm finds the maximum WordNet path similarity with any synset in $S_2$. The same procedure is performed in the opposite direction to ensure symmetry.

The similarity between two descriptions is defined as:

\begin{equation}
Sim(S_1,S_2)=
\frac{
\sum_{s\in S_1} \max_{s'\in S_2} path(s,s') +
\sum_{s'\in S_2} \max_{s\in S_1} path(s',s)
}{2 \cdot \max(|S_1|,|S_2|)}
\end{equation}

where $path(s,s')$ denotes the WordNet path similarity between two synsets.

This procedure produces a $50 \times 50$ similarity matrix representing semantic similarity between all pairs of attack descriptions.

\subsection{Results}

Figures~\ref{fig:WNbased_cfg01}, \ref{fig:WNbased_cfg02}, \ref{fig:WNbased_cfg03}, and \ref{fig:WNbased_cfg04} show the similarity matrices obtained under the four preprocessing configurations.

\begin{figure}[h]
\centering
\includegraphics[width=0.75\linewidth]{step3/wordnet_cfg01_withSW_noLem.png}
\caption{WordNet-based similarity matrix for configuration cfg01 (stopwords included, no lemmatization).}
\label{fig:WNbased_cfg01}
\end{figure}

\begin{figure}[h]
\centering
\includegraphics[width=0.75\linewidth]{step3/wordnet_cfg02_withSW_lem.png}
\caption{WordNet-based similarity matrix for configuration cfg02 (stopwords included, lemmatization applied).}
\label{fig:WNbased_cfg02}
\end{figure}

\begin{figure}[h]
\centering
\includegraphics[width=0.75\linewidth]{step3/wordnet_cfg03_noSW_noLem.png}
\caption{WordNet-based similarity matrix for configuration cfg03 (stopwords removed, no lemmatization).}
\label{fig:WNbased_cfg03}
\end{figure}

\begin{figure}[h]
\centering
\includegraphics[width=0.75\linewidth]{step3/wordnet_cfg04_noSW_lem.png}
\caption{WordNet-based similarity matrix for configuration cfg04 (stopwords removed and lemmatization applied).}
\label{fig:WNbased_cfg04}
\end{figure}

As shown in Figures~\ref{fig:WNbased_cfg01}--\ref{fig:WNbased_cfg04}, all matrices contain a strong diagonal with similarity values equal to 1, representing identical document comparisons. Off-diagonal values represent semantic similarity between different cyber-attack consequence descriptions. \\
Configuration cfg01 (Figure~\ref{fig:WNbased_cfg01}) yields lower similarity values due to the absence of lexical normalization. 
Applying lemmatization in cfg02 (Figure~\ref{fig:WNbased_cfg02}) increases synset coverage and improves semantic matching. 
Configuration cfg03 (Figure~\ref{fig:WNbased_cfg03}) improves similarity detection by removing stopwords and focusing on content-bearing terms.
Finally, cfg04 (Figure~\ref{fig:WNbased_cfg04}) combines stopword removal with lemmatization and produces the most consistent similarity patterns. This configuration reduces noise while improving WordNet synset matching.\\


\section{Matrices Compare}
\subsection{Methods}

\subsection{Results}


\begin{figure}[h]
\centering
\includegraphics[width=1.1\linewidth]{step4/comparison_bar_chart.png}
\caption{NLP Vs. Human}
\label{fig:comparison}
\end{figure}


\section{Conclusion}


\bibliographystyle{IEEEtran}
\bibliography{references}


\end{document}
